% !TEX root = ../libretto.tex
% Sezione: Altri canti
\momento{Altri canti}
\begin{songs}{titoli}

% Fonte: testo fornito dall'utente (tolti gli accordi; corretti i caratteri rovinati: tutti, restiamo, gettiamo, attimo).
% Nell'elenco: "Beato il cuore che perdona". Inno della GMG di Cracovia 2016 (versione italiana).
\beginsong{Beato il cuore che perdona}[by={J. Blycharz}]
\beginverse
Sei sceso dalla tua immensità
in nostro aiuto.
Misericordia scorre da te
sopra tutti noi.
Persi in un mondo d'oscurità
lì tu ci trovi.
Nelle tue braccia ci stringi e poi
dai la vita per noi.
\endverse
\beginchorus
Beato è il cuore che perdona!
Misericordia riceverà da Dio in cielo!
\endchorus
\beginverse
Solo il perdono riporterà
pace nel mondo.
Solo il perdono ci svelerà
come figli tuoi.
\endverse
\beginverse
Col sangue in croce hai pagato tu
le nostre povertà.
Se noi ci amiamo e restiamo in te
il mondo crederà!
\endverse
\beginverse*
Le nostre angosce ed ansietà
gettiamo ogni attimo in te.
Amore che non abbandona mai,
vivi in mezzo a noi!
\endverse
\endsong

% Fonte: testo fornito dall'utente (tolti i vocalizzi "oh oh"). Autore non indicato.
\beginsong{Cerco il tuo volto}
\beginverse
Contemplo in silenzio
un quadro scomposto,
ricerco il mio volto,
desidero te.
Mi incontri spandendo,
sprecando l'amore,
consegni la vita
umiliato per me.
Accetta, Signore, le lacrime amare
di chi ha rinnegato e venduto l'amore.
\endverse
\beginchorus
Io cerco il tuo volto,
il tuo sguardo d'amore
per chi ha rinnegato
e venduto l'amore.
Cerco il tuo volto, cerco il tuo volto,
cerco il tuo volto,
il tuo volto, il tuo volto, Signore.
\endchorus
\beginverse
Tra polvere e rabbia
che ingombra il mio cuore
riveli il tuo volto
di Figlio che muore.
Quando un gelido masso
serra forte l'amore,
io ti riconosco
Maestro e Signore.
Accetta, Signore, il grido stupito
di chi t'ha incontrato e ti vuole annunciare.
\endverse
\beginverse
Ho trovato il tuo volto,
il tuo sguardo soave,
la segreta dolcezza
che mi vuol perdonare.
Ho trovato il tuo volto,
il tuo sguardo soave:
ora vado nel mondo
ed insegno a pregare.
\endverse
\beginchorus
Cerco il tuo volto, cerco il tuo volto,
cerco il tuo volto,
il tuo volto, il tuo volto, Signore.
\endchorus
\endsong

% Fonte: "Con-Te-Camminero-Official-Lead-Sheet.pdf" (melodia con accordi).
% Nell'elenco: senza tipologia/periodo.
\beginsong{Con te camminerò}[by={Gen Verde}]
\beginverse
Mi hai chiamato, ora sono qui.
La mia vita ormai, Signore, trova senso solo in te.
Cerco nel profondo e tu ci sei.
Guardo intorno e poi ti trovo accanto a me.
\endverse
\beginchorus
Insieme a te camminerò.
In ogni volto, in ogni pianto, io ti riconoscerò.
Ti seguirò ovunque andrai.
Io con te camminerò.
\endchorus
\beginverse
Come servo sei venuto qui,
per chi è solo e chi è ferito, per chi vive in povertà.
Ora, o Signore, manda me,
segno del tuo amore per l'umanità.
\endverse
\beginverse*
Tu sei qui, sei con noi, sempre accanto a noi.
Nel dolore di chi non ha più una casa, tu sei lì.
Nel silenzio di chi ha perso ogni speranza, tu sei lì.
Nei rimpianti, nei miei dubbi
e nel buio delle mie fragilità, tu sei lì.
\endverse
\beginverse*
And I will walk on holy ground.
I'll hear your voice, I'll see your face, I will find you in the crowd.
Lord, I will go where you will lead me.
I will walk on holy ground.
\endverse
\beginverse*
Et avec toi, j'irai, Seigneur.
J'entends ta voix, et je te vois, dans les pauvres autour de moi.
Je te suivrai où tu me conduis.
Je te vois auprès de moi.
\endverse
\beginverse*
You are near, you are here.
Io con te camminerò.
\endverse
\endsong

% Fonte: LibrettoAccordi.pdf (Canti Coro), F1.
% Nell'elenco: Credo, Ordinario. In Altri canti su indicazione dell'utente (unico Credo, sezione Credo eliminata).
\beginsong{Credo in unum Deum}[by={M. Balduzzi (CD Verbum Panis)}]
\beginverse
Credo in unum Deum Patrem omnipotentem.
Credo in unum Deum factorem coeli et terrae,
visibilium omnium et invisibilium
Credo in unum Deum. Amen.
\endverse
\beginverse
Credo in unum Deum Dominum Jesum Christum
Credo in unum Deum Filium Dei unigenitum
et ex Patre natum ante omnia saecula
Credo in unum Deum. Amen.
\endverse
\beginverse
Credo in unum Deum Spiritum Sanctum
Credo in unum Deum Dominum et vivificantem
qui ex Patre Filioque procedit
Credo in unum Deum. Amen.
\endverse
\beginverse
Credo in unum Deum Patrem omnipotentem
Credo in unum Deum Dominum Jesum Christum
Credo in unum Deum Spiritum Sanctum
Credo in unum Deum. Amen. Amen.
\endverse
\endsong

% Fonte: testo fornito dall'utente.
% Nell'elenco: "Incendio", Fraternità & Serena Lambertucci, Adorazione; l'utente lo mette in Altri canti.
\beginsong{Incendio}[by={Fraternità, S. Lambertucci}]
\beginverse
Adesso come fuoco partiremo per il mondo
e canteremo la parola d'ogni uomo:
fraternità.
Eccoci, manda noi come luce nell'oscurità,
porteremo salvezza in ogni estremità
della terra.
\endverse
\beginchorus
Tu ci hai scelto dal mondo perché portiamo frutto,
come fosse un incendio:
tu sarai insieme a noi in ogni attimo per farci tuoi. \rep{2}
Per farci tuoi.
\endchorus
\beginverse
Adesso come membra del tuo corpo cureremo
ogni cuore che ha bisogno di conforto:
fraternità.
L'amore che tu ci hai donato ci guiderà,
non sarà più da sola ogni singola anima
in eterno.
\endverse
\beginverse*
Non temiamo perché tu ci hai chiamato
e lo Spirito arde in noi. \rep{4}
Arde in noi.
\endverse
\beginchorus
Tu ci hai scelto dal mondo perché portiamo frutto,
come fosse un incendio:
tu sarai insieme a noi in ogni attimo per farci tuoi. \rep{3}
Per farci tuoi.
\endchorus
\endsong

% Fonte: LibrettoAccordi.pdf (Canti Coro), J17.
% Nell'elenco: senza tipologia (nel libretto accordi tra i canti di comunione). L'utente lo mette in Altri canti.
\beginsong{Nascerà}[by={Gen Rosso}]
\beginverse
Non c'è al mondo chi mi ami, non c'è stato mai nessuno,
in fondo alla mia vita, come Te.
È con te la mia partita, come sabbia tra le dita
scorrono i miei giorni insieme a Te.
Inquietudine o malinconia:
non c'è posto per loro in casa mia.
Sempre nuovo è il tuo modo di
inventare il gioco del tempo per me.
\endverse
\beginchorus
Nascerà dentro me, sul silenzio che abita qui.
Fiorirà un canto che mai nessuno ha cantato per Te.
\endchorus
\beginverse
Se la strada si fa dura come posso aver paura?
Nel buio della notte ci sei Tu.
Se mi assale la fatica di cancellare la sconfitta,
dietro ogni ferita sei ancora Tu.
È una cosa che non mi spiego mai,
cosa ho fatto perché tu scegliessi me?
Cosa mai dirò quando mi vedrai,
quando dai confini del mondo verrai?
\endverse
\endsong

% Fonte: LibrettoAccordi.pdf (Canti Coro), I3 "Segno della pace".
% Nell'elenco: "Pace", Balduzzi, Pace, Ordinario. CD: Giorno senza tramonto (2013).
% Omessa la ripresa finale a due cori (le due parti sovrapposte).
\beginsong{Pace}[by={M. Balduzzi (CD Giorno senza tramonto)}]
\beginchorus
Pace! Pace! Segno della pace che hai portato agli uomini.
Pace! Pace! Pace che qui in terra hai dato a tutti gli uomini.
\endchorus
\beginverse*
Dove c'è odio porterò pace, dov'è tempesta il sole,
dov'è silenzio sarò una voce, dov'è deserto una fonte.
In mezzo al gelo porterò fuoco, in mezzo a spine una rosa,
dove c'è noia porterò vita, dov'è tristezza un sorriso.
\endverse
\endsong

\end{songs}
